\chapter{Free fields in two holomorphic dimensions}
\label{chap:two-dimensional-free-fields}

The residue module records which singular coefficients can occur
near a collision. A field theory must also determine their
coefficients from its kinetic operator. For a free theory this is
an inverse problem: construct a homotopy for the field differential,
then contract observables with its kernel.

The Gaussian calculations of
Chapter~\ref{chap:gaussian-observables} supply the algebra of
contractions. Spatial dependence changes the finite covariance
matrix into a differential kernel. On $\mathbb C^2$, that kernel
is a one-form on the punctured space. Its residue pairs the
two-variable jets constructed in
Proposition~\ref{prop:ocr-residue-dual}.
This pairing gives a Heisenberg algebra of free modes.
A separate central three-input operation gives the abelian higher
current algebra. We construct both, with their different degrees.

\section{Fields and the holomorphic differential}

Write $z_j=x_j+iy_j$, for $j=1,2$, and fix
\[
 \Omega=dz_1\wedge dz_2,\qquad
 dV_4=dx_1\wedge dy_1\wedge dx_2\wedge dy_2.
\]
A smooth $(0,q)$-form is a sum of smooth functions multiplied by
wedge products of $q$ of the symbols $d\bar z_1,d\bar z_2$.
Write their space as $\Omega^{0,q}$.
The wedge symbols anticommute.
The differential
\[
 \bar\partial=\sum_{j=1}^2d\bar z_j\wedge\partial_{\bar z_j},
 \qquad
 \partial_{\bar z_j}=\tfrac12(\partial_{x_j}+i\partial_{y_j}),
 \quad
 \partial_{z_j}=\tfrac12(\partial_{x_j}-i\partial_{y_j}),
\]
squares to zero: the scalar derivatives commute, whereas the wedge
symbols anticommute.
A smooth function killed by $\bar\partial$ is holomorphic.

Let $V$ be a finite-dimensional complex vector space and let
$V^\vee=\operatorname{Hom}_{\mathbb C}(V,\mathbb C)$ be its dual.
The cotangent pair consists of fields
\[
 \gamma\in\Omega^{0,\bullet}\otimes V,
 \qquad
 \beta\in\Omega^{2,\bullet}\otimes V^\vee[1].
\]
Here $\Omega^{2,q}=\Omega^{0,q}\wedge\Omega$.
The shift is on the coefficient space.
Thus the field-complex degrees are $q$ for $\gamma^{0,q}$ and
$q-1$ for $\beta^{2,q}$.
A component coordinate on a field space has the negative of that
space's degree. Its ghost number is therefore $-q$ for $\gamma$
and $1-q$ for $\beta$.
Total degree in products of superfields means ghost number plus
the antiholomorphic form degree.

Take compact support whenever an integral uses both fields.
The action and the odd pairing are
\[
 S_{\mathrm{free}}=\int_{\mathbb C^2}
             \langle\beta,\bar\partial\gamma\rangle,
 \qquad
 \omega_{\mathrm{free}}=\int_{\mathbb C^2}
             \langle\delta\beta,\delta\gamma\rangle.
\]
The angle brackets evaluate a covector on a vector, with the
graded wedge signs. The variation symbol $\delta$ here varies
fields; it is distinct from the spatial differential $\bar\partial$.
The action selects $q_\beta+q_\gamma+1=2$.
Its ghost number is $1-q_\beta-q_\gamma=0$.
The pairing selects $q_\beta+q_\gamma=2$ and has degree $-1$.

Integration by parts has no boundary term on compact supports.
Varying the action therefore gives the linear field differential
$Q=\bar\partial$ through the displayed odd pairing.
The resulting classical master equation is the identity $Q^2=0$,
with the pairing invariant under $Q$.
All these statements concern the quadratic functional and its
formal Batalin--Vilkovisky operations.

\section{A Green homotopy with fixed metric conventions}

To invert $Q$ up to homotopy, choose the coefficient inner product
in which $d\bar z_1,d\bar z_2$ are orthogonal with squared norm
$4$. Contraction $\iota_j$ deletes $d\bar z_j$ from a wedge word,
with the sign of its position, and gives zero if it is absent.
It satisfies
\[
 (d\bar z_i\wedge)\iota_j+\iota_j(d\bar z_i\wedge)
 =\delta_{ij}.
\]
The formal adjoint of $Q$ is then
\[
 Q^{\mathrm{GF}}=-4\sum_j\partial_{z_j}\iota_j.
\]
The superscript denotes this chosen operator of degree $-1$.
Using the contraction identity gives
\begin{equation}\label{eq:c2-laplacian}
 QQ^{\mathrm{GF}}+Q^{\mathrm{GF}}Q
 =-4\sum_j\partial_{z_j}\partial_{\bar z_j}
 =-\Delta_{\mathbb R^4},
\end{equation}
where $\Delta_{\mathbb R^4}$ is the sum of the four real second
derivatives. Its negative is the positive Laplace operator.

Let $r=|z|$ and define
\[
 G_4(z)=\frac1{4\pi^2r^2}.
\]
This function is locally integrable at zero: the radial volume
element contains $r^3dr$, and $r^3r^{-2}$ is integrable.
For a radial function $g(r)$, differentiating in Cartesian
coordinates gives
$\Delta g=g''+3r^{-1}g'$ away from zero.
Thus $\Delta G_4=0$ there.
The outward normal flux of $-\nabla G_4$ through the sphere of
radius $r$ is
\[
 \frac1{2\pi^2r^3}\operatorname{area}(S^3_r)=1,
 \qquad \operatorname{area}(S^3_r)=2\pi^2r^3.
\]
The area follows, for example, by evaluating the four-dimensional
Gaussian integral in polar coordinates.
Integration by parts outside a ball of radius $\epsilon$, followed
by $\epsilon\to0$, therefore gives
$(-\Delta)G_4=\delta_0$.
Here $\delta_0$ denotes the distribution sending a smooth test
function to its value at zero.

For a compactly supported smooth form $f$, define convolution
coefficientwise by
$(G_4*f)(z)=\int G_4(z-w)f(w)\,dV_4(w)$.
Distributional differentiation commutes with this convolution.
Consequently
\begin{equation}\label{eq:c2-green-homotopy}
 H=Q^{\mathrm{GF}}(G_4*-),\qquad QH+HQ=1.
\end{equation}
The output need not have compact support.
Equation~\eqref{eq:c2-green-homotopy} is an identity on the full
Euclidean space with compactly supported inputs.

The positive-time heat kernel is
\[
 k_\tau(z)=(4\pi\tau)^{-2}
               \exp\bigl(-|z|^2/(4\tau)\bigr),\qquad\tau>0.
\]
The Gaussian integral gives its integral one.
Differentiation verifies $\partial_\tau k_\tau=\Delta k_\tau$.
Substitution $u=r^2/(4\tau)$ in its integral gives
$\int_0^\infty k_\tau(z)d\tau=G_4(z)$ for $z\ne0$.
The finite interval $\epsilon\le\tau\le L$ therefore gives a
smooth homotopy kernel with differential $k_\epsilon-k_L$.

\section{The normalized two-dimensional contraction}

Fix the outward orientation on every sphere and put
\begin{equation}\label{eq:c2-bm-kernel}
 B(z)=\frac1{4\pi^2}
 \frac{\bar z_1\,d\bar z_2-\bar z_2\,d\bar z_1}{|z|^4},
 \qquad K(z)=B(z)\wedge\Omega.
\end{equation}
The numerator has one antiholomorphic differential.
Thus $B$ has Dolbeault degree one.

\begin{proposition}\label{prop:c2-normalized-kernel}
The kernel satisfies
\[
 \bar\partial B=0\quad(z\ne0),\qquad
 \int_{S^3_\epsilon}K=1,
 \qquad \bar\partial K=\delta_0\,dV_4
\]
in the last equation as a current.
\end{proposition}

\begin{proof}
Set $s=|z|^2$ and
$A=\bar z_1d\bar z_2-\bar z_2d\bar z_1$.
Then
\[
 \bar\partial A=2d\bar z_1\wedge d\bar z_2,
 \qquad
 \bar\partial s\wedge A=s\,d\bar z_1\wedge d\bar z_2.
\]
The two terms in $\bar\partial(s^{-2}A)$ cancel.

On the sphere, $A\wedge\Omega$ is invariant under unitary
coordinate changes. Such a change multiplies $A$ by the conjugate
determinant and $\Omega$ by the determinant, whose product is one.
At $(\epsilon,0)$ it equals
$2\epsilon\,dy_1\wedge dx_2\wedge dy_2$.
This is $2\epsilon$ times the outward surface measure there,
and hence everywhere by unitary invariance.
The area formula gives
$(4\pi^2\epsilon^4)^{-1}(2\epsilon)(2\pi^2\epsilon^3)=1$.

The coefficients of $K$ have size $O(r^{-3})$, so are locally
integrable in four dimensions.
For a test function, apply Stokes' formula outside the ball of
radius $\epsilon$.
The integral over its boundary tends to the value at zero.
Indeed, replacing the test function by that value has error
$O(\epsilon)$ times a bounded angular integral.
Since $\partial K=0$ by its holomorphic top degree, this proves
the stated $\bar\partial$ current identity.
\end{proof}

The orientation matters: $1/(2\pi i)^2=-1/(4\pi^2)$ would give
mass $-1$ for the same numerator and outward sphere.
With the convention~\eqref{eq:c2-bm-kernel}, the heat regularization is
\[
 B_{\epsilon,\infty}(z)
 =B(z)\left[1-\left(1+\frac{|z|^2}{4\epsilon}\right)
                   e^{-|z|^2/(4\epsilon)}\right].
\]
Differentiating the bracketed radial factor and using
$d\bar z_1\wedge d\bar z_2\wedge\Omega=4dV_4$ gives
\[
 \bar\partial(B_{\epsilon,\infty}\wedge\Omega)
 =k_\epsilon\,dV_4.
\]
The difference
$B_{\epsilon,L}=B_{\epsilon,\infty}-B_{L,\infty}$
therefore has differential $(k_\epsilon-k_L)dV_4$.

For two points, the full kernel uses
$d\overline{z_j-w_j}=d\bar z_j-d\bar w_j$.
The ordered free contraction is normalized by
\[
 \beta_\lambda(z)\gamma_v(w)
 \big|_{\mathrm{contraction}}
 =\hbar\,\lambda(v)B(z-w),\qquad
 v\in V,\quad\lambda\in V^\vee.
\]
This notation denotes the single paired term in the Wick expansion.
The coefficient $\hbar$ is the formal loop parameter.
Difference forms and the graded pairing determine the reversed
contraction and the descendant components.
The form degree one combines with the coefficient-field shifts
to give the propagator total degree zero.

The kernel is smooth at every nonzero point of either coordinate
axis. Its pullback to the punctured line $z_2=0$ is zero, because
both $\bar z_2$ and the pullback of $d\bar z_2$ vanish.
A nonzero one-dimensional singular coefficient therefore requires
a transverse residue operation such as the derived restriction
constructed in Chapter~\ref{chap:ordered-collisions-residues}.

\section{Residue derivatives and two-variable modes}

For nonnegative integers $a,b$, define
\[
 B_{ab}=\frac{(-1)^{a+b}}{a!b!}
              \partial_{z_1}^a\partial_{z_2}^bB.
\]
For a holomorphic function $f$ near a closed ball containing zero,
Stokes' formula and Proposition~\ref{prop:c2-normalized-kernel} give
$\int_{\partial D}f(z)K(z-w)=f(w)$ for $w$ in its interior.
Differentiate this identity with respect to $w$ while keeping the
boundary fixed. The kernel is smooth there, so differentiation
passes under the integral. Since $\partial_wB(z-w)=-\partial_zB(z-w)$,
\[
 \int_{\partial D} f B_{ab}\wedge\Omega
 =\frac1{a!b!}\partial_{z_1}^a\partial_{z_2}^bf(0).
\]
This is exactly the Taylor coefficient functional paired with
$\rho_{ab}$ in Proposition~\ref{prop:ocr-residue-dual}.
The smooth kernel on the puncture and the algebraic principal part
therefore give the same normalized residue functional.
The algebraic class has degree one on the puncture and degree two
in the support complex.

Let $A_0=\mathbb C[z_1,z_2]$ and let
$P=\bigoplus_{a,b\ge0}\mathbb C e_{ab}$, where $e_{ab}$ denotes
the principal part $[z_1^{-a-1}z_2^{-b-1}]$ without its volume symbol.
Use the pairing
\[
 \langle e_{ab},z_1^mz_2^n\rangle
 =\delta_{am}\delta_{bn}.
\]
Assign degree zero to $A_0$ and degree one to $P$.
No arbitrary infinite sum of principal parts is allowed.

\section{The free Heisenberg algebra and its representation}

For a graded vector space, a Lie bracket of degree zero obeys
$[u,v]=-(-1)^{|u||v|}[v,u]$ and the graded Jacobi identity
\[
 [u,[v,w]]=[[u,v],w]+(-1)^{|u||v|}[v,[u,w]].
\]
Introduce symbols $c(a,v)$ and $b(a,\lambda)$, linear in their
arguments, with $a\in A_0\oplus P$.
Their degrees are
\[
 \begin{array}{c|cccc}
 \text{symbol}&c(f,v)&c(u,v)&b(f,\lambda)&b(u,\lambda)\\
 \text{degree}&0&1&-1&0
 \end{array}
 \qquad(f\in A_0,\ u\in P).
\]
Add a degree-zero symbol $K$.
An element is central when its bracket with every element is zero.
Declare $K$ central, set the brackets of two $b$'s or two $c$'s
to zero, and put
\begin{equation}\label{eq:c2-heisenberg-bracket}
 \begin{split}
 [b(u,\lambda),c(f,v)]&=\lambda(v)\langle u,f\rangle K,\\
 [b(f,\lambda),c(u,v)]&=\lambda(v)\langle u,f\rangle K.
 \end{split}
\end{equation}
All other mixed brackets are zero, and reversed orders are given
by graded antisymmetry.
Every nonzero bracket has degree zero and central output.
Every nested bracket therefore vanishes, which proves Jacobi.
These formulas define the free complex-two-dimensional Heisenberg
mode algebra $\mathsf H_{2,V}$.

For $V=\mathbb C$, its even oscillator relation is
\[
 [b(e_{ab}),c(z_1^mz_2^n)]
 =\delta_{am}\delta_{bn}K.
\]
There is one conjugate pair for each two-variable Taylor coefficient.
The odd pair in~\eqref{eq:c2-heisenberg-bracket} obeys an
anticommutator, since its two degrees are odd.

Both signs can be realized directly.
Form the polynomial algebra on even generators indexed by
$A_0\otimes V$, and tensor it with the exterior algebra on
degree-minus-one generators indexed by $A_0\otimes V^\vee$.
An element involves finitely many such generators.
Let $c(f,v)$ act by even multiplication and $b(f,\lambda)$ by
odd multiplication. Let $b(u,\lambda)$ act by the corresponding
even residue derivative, multiplied by $\hbar$.
Let $c(u,v)$ act by the odd residue derivative, also multiplied by
$\hbar$.
An odd derivative deletes the corresponding exterior generator
with its position sign.

The ordinary derivative-multiplication commutator is evaluation.
The odd derivative-multiplication anticommutator is evaluation,
because moving the deleted generator through an exterior word
cancels every other term.
Thus the representation satisfies~\eqref{eq:c2-heisenberg-bracket}
with $K$ acting as $\hbar$.
Each operator is defined on every polynomial state by a finite sum.

\section{A central operation with three inputs}

The free oscillator pairing is bilinear.
Abelian currents on a holomorphic surface also admit a central
operation involving one singular coefficient and two functions.
The two derivatives of those functions supply the holomorphic
volume degree required by the residue.

Put $A=\mathbb C[[z_1,z_2]]$ and define
\[
 \{f,g\}=\partial_{z_1}f\,\partial_{z_2}g
          -\partial_{z_2}f\,\partial_{z_1}g.
\]
The derivatives and products are formal and coefficientwise.
A principal part tests only finitely many coefficients of this
expression, so its residue is defined for all $f,g\in A$.

Use the graded space with $A\oplus\mathbb CK$ in degree zero
and $P$ in degree one.
An $L_\infty$ algebra has graded antisymmetric operations
$\ell_n$ of degree $2-n$.
To state their identities, let $\operatorname{Sh}(i,n-i)$ consist
of permutations $\sigma$ for which
$\sigma(1)<\cdots<\sigma(i)$ and
$\sigma(i+1)<\cdots<\sigma(n)$.
For homogeneous inputs $v_1,\ldots,v_n$, let $\chi(\sigma;v)$
be the sign obtained by exchanging adjacent inputs using the
factor $-(-1)^{|v_a||v_b|}$ for each exchange.
The generalized Jacobi identity, for each $n\ge1$, is
\[
 \sum_{i+j=n+1}(-1)^{i(j-1)}
 \sum_{\sigma\in\operatorname{Sh}(i,n-i)}\chi(\sigma;v)
 \ell_j\bigl(\ell_i(v_{\sigma(1)},\ldots,v_{\sigma(i)}),
             v_{\sigma(i+1)},\ldots,v_{\sigma(n)}\bigr)=0.
\]
In particular, its first operation $\ell_1$ is a differential.
For a family with only a central $\ell_3$, every such identity
vanishes term by term: every possible inner operation is either
zero or central, and a central input kills the outer operation.
This supplies all identities for the following explicit family.

Fix $\kappa\in\mathbb C$ and define
\begin{equation}\label{eq:c2-ternary-residue}
 \ell_3(u,f,g)=\kappa\operatorname{Res}
                    \bigl(u\{f,g\}\Omega\bigr)K,
 \qquad u\in P,\quad f,g\in A.
\end{equation}
Extend by graded antisymmetry.
Set every other input degree pattern, every central input, and
every $\ell_n$ with $n\ne3$ equal to zero.
The input degree in~\eqref{eq:c2-ternary-residue} is one and its
output degree is zero, so $\ell_3$ has the required degree $-1$.
The central-output observation proves all higher identities.

For monomials the operation is
\begin{equation}\label{eq:c2-ternary-modes}
 \ell_3(e_{ab},z_1^mz_2^n,z_1^pz_2^q)
 =\kappa(mq-np)
       \delta_{a,m+p-1}\delta_{b,n+q-1}K.
\end{equation}
Indeed, the determinant of derivatives is
$(mq-np)z_1^{m+p-1}z_2^{n+q-1}$, and the residue extracts
exactly the two indicated exponents.
Terms requiring a negative index vanish.
In particular,
\[
 \ell_3(e_{00},z_1,z_2)=\kappa K,
 \qquad
 \ell_3(e_{10},z_1^2,z_2)=2\kappa K.
\]
The latter survives even if the singular coefficient module is
restricted to the annihilator of constant functions.

This abelian higher current algebra has zero binary bracket.
Its parameter $\kappa$ specifies a central three-input operation.
Determining $\kappa$ from a quantum field theory requires a
comparison with its current insertions and their three-point
contractions. The free residue pairing alone does not specify
that comparison.

\section{Adding the two topological directions}

The mixed local geometry is
\[
 X=\mathbb R_t\times\mathbb R_s\times\mathbb C^2_{z_1,z_2}.
\]
Its linear differential treats the real and holomorphic directions
differently:
\[
 Q=dt\wedge\partial_t+ds\wedge\partial_s+\bar\partial.
\]
All four one-form symbols anticommute, so $Q^2=0$.
The same contraction rule as above gives
\[
 Q\iota_{dt}+\iota_{dt}Q=\partial_t,\qquad
 Q\iota_{ds}+\iota_{ds}Q=\partial_s,\qquad
 Q\iota_{d\bar z_j}+\iota_{d\bar z_j}Q=\partial_{\bar z_j}.
\]
Thus real and antiholomorphic translations have explicit
$Q$-homotopies. A holomorphic derivative need not act trivially
on cohomology: $\partial_{z_1}$ sends the holomorphic function
$z_1$ to $1$, and neither is a boundary in degree zero.
This distinguishes the holomorphic directions in the mixed complex.
The field coefficients must depend jointly and smoothly on all
six real coordinates. Define the mixed form complex directly by
\[
 \mathcal M(X)=C^\infty(X;\mathbb C)
      \otimes_{\mathbb C}\Lambda(dt,ds,d\bar z_1,d\bar z_2),
\]
with differential $Q$ as displayed above.
The exterior factor is finite-dimensional, so this tensor product
contains every jointly smooth mixed form.
For a finite-dimensional coefficient space $W$, take
\[
 \alpha\in\mathcal M(X)\otimes W[1],\qquad
 \beta\in\mathcal M(X)\otimes W^\vee[2].
\]
A component of total form degree $r$ has field-complex degree
$r-1$ in $\alpha$ and $r-2$ in $\beta$.
On fields of degree zero, $\alpha$ is a mixed one-form and
$\beta$ is a mixed two-form.
The field strength is $F=Q\alpha$.
Pairing it with the multiplier $\beta$ gives the free BF action
\[
 S_{\mathrm{mix}}=\int_X\langle\beta,Q\alpha\rangle\wedge\Omega.
\]
Varying $\beta$ imposes the equation $F=0$.
For a smooth $W$-valued function $\varepsilon$, the change
$\alpha\mapsto\alpha+Q\varepsilon$ preserves the action because
$Q^2\varepsilon=0$.
For a mixed $W^\vee$-valued one-form $\eta$, the change
$\beta\mapsto\beta+Q\eta$ also preserves it:
the change is the integral of
$Q\langle\eta,Q\alpha\rangle\wedge\Omega$, which vanishes
on compact supports. These are the gauge symmetries of the
free mixed action.
The integral selects total mixed form degree four.
The associated pairing has degree $-1$, and the action has degree
zero by the same component count as in the holomorphic theory.
This specifies the free cotangent theory with finite coefficients.

Give $dt,ds$ squared norm one and $d\bar z_j$ squared norm four.
With contractions defined by deleting these symbols, put
\[
 Q^{\mathrm{GF}}=-\partial_t\iota_{dt}-\partial_s\iota_{ds}
                  -4\sum_j\partial_{z_j}\iota_{d\bar z_j}.
\]
The contraction identities give
$QQ^{\mathrm{GF}}+Q^{\mathrm{GF}}Q=-\Delta_{\mathbb R^6}$.
Write $R^2=t^2+s^2+|z_1|^2+|z_2|^2$.
The scalar Green function and heat kernel are
\[
 G_6=\frac1{4\pi^3R^4},\qquad
 k_\tau^{\mathrm{mix}}=(4\pi\tau)^{-3}e^{-R^2/(4\tau)}.
\]
The radial Laplacian is $g''+5R^{-1}g'$.
The sphere area is $\operatorname{area}(S^5_R)=\pi^3R^5$.
Thus the outward flux of $-\nabla G_6$ is one, proving
$(-\Delta)G_6=\delta_0$ as before.
It follows that $H_{\mathrm{mix}}=Q^{\mathrm{GF}}(G_6*-)$ satisfies
$QH_{\mathrm{mix}}+H_{\mathrm{mix}}Q=1$ on compactly supported
smooth inputs.

The corresponding normalized fundamental current is explicit.
Set
\[
 \nu=dt\wedge ds\wedge d\bar z_1\wedge d\bar z_2,
 \qquad
 W_0=t\partial_t+s\partial_s
             +2\bar z_1\partial_{\bar z_1}
             +2\bar z_2\partial_{\bar z_2}.
\]
Then
\begin{equation}\label{eq:c2-mixed-kernel}
 K_{\mathrm{mix}}=
       \frac{\iota_{W_0}\nu\wedge\Omega}{4\pi^3R^6},
 \qquad QK_{\mathrm{mix}}=\delta_0\,dV_6,
\end{equation}
where $dV_6=dt\wedge ds\wedge dV_4$.
Away from zero, differentiating its coefficient gives
\[
 \partial_t\frac t{R^6}+\partial_s\frac s{R^6}
 +\sum_{j=1}^2\partial_{\bar z_j}
                    \frac{2\bar z_j}{R^6}
 =\frac6{R^6}-\frac{6R^2}{R^8}=0.
\]
The delta normalization is the Green normalization above with
$\nu\wedge\Omega=4dV_6$.
The kernel has size $O(R^{-5})$ and is locally integrable.
Its mixed form degree is three. The two coefficient shifts sum
to $-3$, so the resulting free propagator has total degree zero.

Replacing the scalar Green integral by
$\int_\epsilon^L k_\tau^{\mathrm{mix}}d\tau$ gives a smooth
finite-scale propagator.
Its differential is the difference of the endpoint heat kernels.
Contracting polynomial observables with it is the spatial form
of the Gaussian homotopy: every contraction removes two input
fields, so the exponential terminates at fixed polynomial degree.

The free kernel determines contractions for the finite-dimensional
coefficient space $W$ on $X$.
An interacting Hamiltonian theory requires its nonlinear field
equations, a coefficient Lie bracket, and a pairing on the coefficient
spaces used in its action and contractions. For infinite-dimensional coefficients, the domain and
topology of the pairing must be specified, and its contractions must
be proved to be defined on the chosen fields.

Write $\mathcal C=A\oplus P\oplus\mathbb CK$ for the graded space
with operations $\ell_n$ constructed above.
A possible target for a current map is a graded algebra $\mathcal B$
with differential $d_{\mathcal B}$ and graded antisymmetric multilinear operations
$b_n:\mathcal B^{\otimes n}\to\mathcal B$ of degree $2-n$,
where $b_1=d_{\mathcal B}$ and the $b_n$ satisfy the generalized
Jacobi identities stated above.
One sufficient form of comparison is a degree-zero linear map
$F:\mathcal C\to\mathcal B$ satisfying, for every $n\ge1$,
\[
 F\bigl(\ell_n(c_1,\ldots,c_n)\bigr)
   =b_n\bigl(F(c_1),\ldots,F(c_n)\bigr)
 \qquad(c_1,\ldots,c_n\in\mathcal C).
\]
This is a strict map of the displayed higher current algebras.
Its equation for $n=1$ preserves the differential; its equation for
$n=3$ identifies the normalized residue operation with the specified
three-input operation on $\mathcal B$.
A comparison permitting these equalities only up to homotopy would
instead require the corresponding homotopies and the equations they
must satisfy.

No algebra $\mathcal B$ of boundary observables, and no such comparison,
has been constructed from the mixed kernel.
A physical realization must specify the boundary conditions on the
fields, construct their observable algebra and its operations, and
derive the residue normalization and the parameter $\kappa$ from
those operations. The kinetic homotopy alone provides none of
these additional constructions.
